\documentclass[tikz,border=10pt]{standalone}
\input{../../../docs/tikz/dag_preamble.tex}

\begin{document}
\begin{tikzpicture}[x=1cm,y=1cm,
  every node/.append style={outer sep=4pt,font=\sffamily\small}]

\dagPaperMetadata
  {Redesigning Service Level Agreements: Equity and Efficiency in City Government Operations}
  {Zhi Liu and Nikhil Garg}
  {ACM Conference on Economics and Computation (EC)}
  {2024; current manuscript revision}
  {https://arxiv.org/abs/2410.14825}

\dagPaperLegendRightOfMetadata{
\node[dag_result,dag_template_legend] (legRes) at (0,0)
  {checked result};
\node[dag_model,dag_template_legend] (legModel) at (4,0)
  {source model};
}
\daglegend{(legRes)(legModel)}{Legend}

\begin{dagPaperBody}
\node[font=\sffamily\bfseries\Large,anchor=west] at (-2.7,-0.9)
  {Current five-proposition theory and selected supporting results};

\node[dag_model] (Model) at (0,-3.6) {
  \textbf{Inspection model}\\
  fixed admitted loads; GPS service; all-request delay and non-inspection costs
};
\node[dag_result] (Tail) at (7,-3.6) {
  \textbf{Equation (1): stationary GPS tail}\\
  typical-admitted-request response has tail at most
  $e^{-(C_b\phi_{k,b}-s_{k,b})z}$
};
\node[dag_result] (Fraction) at (14,-3.6) {
  \textbf{All-request SLA consequence}\\
  timely-request intensity fraction is at least
  $(s_{k,b}/\lambda_{k,b})(1-\alpha)$
};

\node[dag_result] (Reform) at (0,-7.8) {
  \textbf{Proposition 2.1}\\
  \textbf{Reciprocal-capacity representation}\\
  original GPS-policy and reduced SLA programs have the same minimizers
};
\node[dag_result] (Efficiency) at (7,-7.8) {
  \textbf{Proposition 2.2}\\
  \textbf{Extreme efficiency}\\
  the square-root allocation gives the unique all-request efficiency minimizer
};
\node[dag_result] (Equity) at (14,-7.8) {
  \textbf{Proposition 2.3: extreme equity}\\
  equal within-category burdens, with an efficiency-best equitable endpoint
};

\node[dag_result] (Nonnegative) at (0,-12.3) {
  \textbf{Relative price of equity}\\
  the efficiency loss from requiring equity is nonnegative
};
\node[dag_result] (City) at (7,-12.3) {
  \textbf{City efficiency endpoint}\\
  pooling categories across Boroughs gives the stated city-wide optimizer
};
\node[dag_result] (Price) at (14,-12.3) {
  \textbf{Proposition 2.4: price of equity}\\
  relative Pearson-divergence identity, upper bound, zero-price characterization,
  and conditional capacity scaling
};

\node[dag_result] (Gain) at (7,-16.8) {
  \textbf{Relative centralization gain}\\
  the gain equals the effective-load-square difference divided by
  excess capacity times the original efficiency cost
};
\node[dag_result] (Centralization) at (14,-16.8) {
  \textbf{Proposition 2.5: centralization}\\
  with at least two Boroughs, every cell's SLA improves strictly;
  the gain--price comparison has the stated Pearson-divergence criterion
};

\draw[dag_dashed_arrow] (Model.east) -- (Tail.west);
\draw[dag_arrow] (Tail.east) -- (Fraction.west);
\draw[dag_dashed_arrow] (Model.south) -- (Reform.north);
\draw[dag_arrow] (Reform.east) -- (Efficiency.west);
\draw[dag_arrow] (Reform.north) -- ++(0,0.65) -| (Equity.north);
\draw[dag_arrow] (Efficiency.south west) -- (Nonnegative.north east);
\draw[dag_arrow] (Equity.south west) -- ++(-0.4,-1.0) -| (Nonnegative.north);
\draw[dag_arrow] (Efficiency.south) -- (City.north);
\draw[dag_arrow] (Equity.south) -- (Price.north);
\draw[dag_arrow] (Efficiency.north east) -- ++(0.5,0.6) -| (Price.east);
\draw[dag_arrow] (City.south) -- (Gain.north);
\draw[dag_arrow] (Gain.east) -- (Centralization.west);
\draw[dag_arrow] (Price.south) -- (Centralization.north);

\end{dagPaperBody}
\end{tikzpicture}
\end{document}
